\chapter{Basic Problems}

% ============ Problem 1 ============ %
% An ordinary problem: question, solution and answer.

\begin{problem}
{
Solve the equation $x^2 - 5x + 6 = 0$.
}
{
We look for two numbers whose sum is $5$ and whose product is $6$: they are $2$
and $3$, so the polynomial factors as
\begin{align}
    x^2 - 5x + 6 &= (x - 2)(x - 3). \label{eq:quadratic-factored}
\end{align}
A product is zero only when one of its factors is zero, so $x - 2 = 0$ or
$x - 3 = 0$. As a check, the quadratic formula gives the same roots:
\begin{align*}
    x &= \frac{5 \pm \sqrt{5^2 - 4 \cdot 6}}{2} = \frac{5 \pm 1}{2}.
\end{align*}
}
{
$x = 2$ or $x = 3$.
}
\end{problem}

% ============ Problem 2 ============ %
% A problem in several parts: the main question has no solution and no answer
% ({}{}), and each subproblem may omit its solution or its answer.

\begin{problem}
{
Let $S_n = 1 + 3 + 5 + \dots + (2n - 1)$ be the sum of the first $n$ odd
numbers.\label{prob:odd-sum}
}{}{}
\end{problem}

\begin{subproblem}
{
Compute $S_1$, $S_2$, $S_3$ and $S_4$, and guess a formula for $S_n$.
\label{prob:odd-sum-guess}
}
{
Adding one odd number at a time,
\begin{align*}
    S_1 &= 1, & S_2 &= 1 + 3 = 4, & S_3 &= 4 + 5 = 9, & S_4 &= 9 + 7 = 16.
\end{align*}
These are the squares of $1$, $2$, $3$ and $4$.
}
{
$S_n = n^2$.
}
\end{subproblem}

\begin{subproblem}
{
Prove the formula of part (a) by induction.
}
{
For $n = 1$, $S_1 = 1 = 1^2$. If $S_n = n^2$ for some $n \ge 1$, then
\begin{align}
    S_{n+1} &= S_n + (2n + 1) = n^2 + 2n + 1 = (n + 1)^2,
    \label{eq:odd-sum-induction}
\end{align}
so the formula also holds for $n + 1$, and by induction for every $n \ge 1$.
This part is a proof, so it has no answer box.
}
{}
\end{subproblem}

\begin{subproblem}
{
Evaluate $S_{10}$.
}
{}
{
$S_{10} = 10^2 = 100$ (no solution is needed, so only the answer is given).
}
\end{subproblem}

\begin{subproblem}
{
How many consecutive odd numbers, starting from $1$, must be added for the sum
to exceed $1000$?
}
{
We need the smallest $n$ with $n^2 > 1000$, that is $n > \sqrt{1000} \approx
31.6$. Indeed $31^2 = 961 \le 1000$ while $32^2 = 1024 > 1000$.
}
{
$32$ odd numbers.
}
\end{subproblem}

% ============ Problem 3 ============ %
% A problem with a standalone TikZ figure (Figures/RightTriangle.tex).

\begin{problem}
{
A right triangle has legs $a = 3$ and $b = 4$, as shown in \figref{fig:triangle}.
Find the length of the hypotenuse $c$ and the angle $\theta$ opposite the
leg $a$.
\begin{figure}[H]
    \centering
    \includegraphics{Figures/RightTriangle.pdf}
    \caption{Right triangle with legs $a$ and $b$.}
    \label{fig:triangle}
\end{figure}
}
{
By the Pythagorean theorem,
\begin{align*}
    c &= \sqrt{a^2 + b^2} = \sqrt{9 + 16} = 5.
\end{align*}
The leg $a$ is opposite $\theta$ and the leg $b$ is adjacent to it, so
\begin{align*}
    \tan\theta &= \frac{a}{b} = \frac{3}{4},
    & \theta &= \arctan\frac{3}{4} \approx 36.87^\circ.
\end{align*}
}
{
$c = 5$ and $\theta = \arctan(3/4) \approx 36.87^\circ$.
}
\end{problem}
